이러한 결과는 유효숫자 교육이 단순 규칙 암기를 넘어 측정 과정 전체에서
`왜 이 자릿수가 필요한가'를 스스로 설명하고 판단할 수 있는 방향으로
재구조화될 필요가 있음을 제시한다.
%
즉, 최소 눈금 단위 해석, 어림 판단, 불확실도의 의미를 실험 상황과 연결하여
체험적으로 학습하도록 하는 설계가 요구된다.
%
또한 계산 단계에서는 평균값의 의미와 불확실도 감소의 원리를 탐구적 맥락에서 다루어야 하며,
필요성 판단에서는 도구의 종류가 아니라 측정의 신뢰도 개념을 중심으로 사고할 수 있도록 안내해야 한다.
%
이를 통해 유효숫자가 단순 표기 규칙이 아니라 과학 탐구에서 신뢰도와 정당성을 확보하는
사고 도구로 정착될 수 있을 것이다.


\hlb{연구의 한계점

실제 도구를 조작하는 측정 상황 아님

보고서 등을 평가해야 함}